Este trabajo cambió el centro del problema respecto del TP1. Allí predominaba la
lógica combinacional: una vez presentes operandos y opcode, el resultado quedaba
determinado. En el TP2 aprendimos a describir un comportamiento que se extiende
en el tiempo. Recibir un byte, esperar el siguiente y responder cuando existe
capacidad son acciones que necesitan memoria, no solo una expresión lógica.

Las máquinas de estado nos permitieron separar estado actual, próximo estado y
salidas. Esa separación volvió explícitas preguntas que una descripción
compacta puede ocultar: qué se conserva mientras no llega un tick, durante qué
ciclo se produce un pulso y a qué estado vuelve el circuito después de un
error. La codificación one-hot facilitó seguir esas decisiones en las trazas y
relacionarlas luego con flip-flops y lógica de decodificación en la síntesis. Al
mismo tiempo, comprobamos que one-hot no reemplaza una rama segura para
combinaciones ilegales.

El receptor mostró que una entrada asíncrona no puede tratarse como un dato ya
alineado con el reloj. El sincronizador reduce la probabilidad de propagar
metastabilidad y el sobremuestreo ubica la observación cerca del centro del bit.
El falso start y el stop inválido hicieron visible otra lección: detectar un
error no alcanza; también hay que decidir cómo recuperar el reposo sin contar
varias veces el mismo nivel defectuoso.

Los FIFO y el backpressure ampliaron esa idea más allá de UART. Aprendimos que
cuando productor y consumidor avanzan a ritmos diferentes, la capacidad y las
reglas de aceptación forman parte de la arquitectura. La profundidad cuatro no
hace más veloz la línea, pero permite desacoplar responsabilidades. De manera
análoga, WAIT\_TX no es una espera accidental: garantiza que el resultado y las
banderas persistan hasta que la escritura pueda aceptarse.

La verificación por capas resultó esencial. Los bancos pequeños permitieron
atribuir fallas a un contador, un orden de bits, un límite de FIFO o una
transición. Recién después la prueba pin a pin comprobó la cooperación de todos
los bloques en 812 transacciones, incluidas 800 pseudoaleatorias reproducibles.
Las formas de onda y los PASS responden la pregunta funcional; la síntesis
explica qué hardware se infirió; la implementación y el timing muestran que ese
hardware puede operar con las restricciones declaradas.

Finalmente, aprendimos a no confundir esos niveles de evidencia. El bitstream,
los 458 endpoints con márgenes positivos y los reportes CDC/DRC cerraron la etapa
de herramientas, pero no demostraron el cable, el conversor USB ni la
observación de los LED. Esos elementos quedaron comprendidos en la validación
sobre la Basys 3, cuyo resultado dependió de las respuestas UART, los registros
de ejecución y los patrones físicos. Esa distinción entre niveles de evidencia,
junto con la secuencia
especificar--modularizar--verificar--implementar--medir, constituye el método que
nos llevamos para diseños futuros con protocolos, estado y temporización.
